\section{Related Work}\label{sec:related-work}

This section reviews prior work along the axes that motivate the design in Sections~\ref{sec:problem}--\ref{sec:implementation}: edge orchestration, constrained execution, secure IoT communication, reconciliation semantics, and reproducible evaluation. The main gap across these areas is the lack of an end-to-end path from Kubernetes-style declarative intent to verifiable execution on constrained embedded devices that cannot act as ordinary cluster nodes.

\subsection{Cloud-Native Edge and IoT Orchestration}
Foundational work on fog and edge computing established the need to move computation, storage, and control closer to data sources in order to reduce latency, improve locality, and support cyber-physical applications~\cite{bonomi_fog,shi_edge,satyanarayanan_edge,chiang_fog}. Kubernetes has become the dominant orchestration substrate for cloud-native services because it provides declarative resource management, controllers, reconciliation loops, and extensibility through Custom Resource Definitions (CRDs) and operators~\cite{kubernetes,kubernetes_operators}. Lightweight distributions such as K3s and MicroK8s reduce deployment overhead and make cluster-style operation feasible at smaller sites~\cite{k3s,microk8s}.

Several systems extend Kubernetes semantics toward the edge. KubeEdge moves parts of Kubernetes control closer to edge nodes, OpenYurt targets cloud-edge application management, Akri discovers and exposes leaf devices to Kubernetes, and Krustlet explored WebAssembly workloads as Kubernetes-scheduled units~\cite{kubeedge,openyurt,akri,krustlet}. Recent Kubernetes-native IoT orchestration based on Stack4Things (S4T), CRDs, and Crossplane abstractions similarly demonstrates the value of declarative resource management for IoT deployments~\cite{dagati_iot_orchestration_s4t}. \revC{That line has since been extended toward a deviceless Function-as-a-Service model, in which users deploy and invoke functions according to capability, location, and load without managing the underlying devices, using Nuclio as the serverless runtime on edge workers and WebSocket tunnelling to reach nodes behind NATs or firewalls~\cite{mirto_deviceless_faas}. That design shares the goal of removing device-level management from the user, but its edge workers are Linux single-board hosts running a container-based function runtime, and the tunnelled node remains the unit of execution; the endpoint considered here executes a WASM module under Zephyr/WAMR with no operating-system userspace, container runtime, or cluster agent, so the mediation boundary sits at the gateway rather than at the device itself.} These approaches motivate the present work, but they generally assume Linux-class edge hosts, gateway nodes, container-compatible execution environments, device plugins, or platform services attached to capable nodes. The target considered here is narrower and more constrained: an embedded endpoint that cannot run a kubelet, container runtime, Kubernetes agent, or full cloud-native network stack.

IoT middleware and industrial edge platforms provide additional context. EdgeX Foundry, FIWARE, oneM2M, OMA Lightweight M2M, and OPC UA address aspects of device modeling, interoperability, telemetry, and industrial communication~\cite{edgex,fiware,onem2m,lwm2m,opcua}. These frameworks are broader than a single orchestration mechanism and often emphasize interoperability or application enablement rather than Kubernetes-native desired-state reconciliation. The contribution is therefore not another general IoT middleware layer, but a focused integration that links cloud-style intent, gateway-mediated secure control, and embedded WASM execution through explicit resource ownership.

\revD{The device twin, or shadow, pattern addresses the same problem from the platform side:
a cloud-resident document holds desired and reported state for a device that cannot be
queried directly, and a managed service reconciles the two as the device reconnects. AWS
IoT Greengrass~\cite{aws_greengrass}, Azure IoT Edge~\cite{azure_iot_edge} and the
LwM2M object model~\cite{lwm2m} have long provided this. The evidence model used here is
the same pattern; what differs is the location of the reconciliation point. In twin-based
platforms authoritative state lives in a proprietary service and reaches cluster tooling
through bridges. Here it is a Kubernetes resource from the outset, so status ownership,
admission control, watches and the controller pattern extend to constrained devices
directly, and a placement policy is an ordinary controller over the same objects.}

\subsection{Constrained Execution and Secure Communication}
Embedded and IoT operating systems provide the substrate on which constrained devices communicate and execute application logic. Zephyr, FreeRTOS, RIOT, and Contiki-NG represent different points in the design space for real-time behavior, connectivity, portability, and resource footprint~\cite{zephyr,freertos,riot,contiki}; TinyOS established the long-standing importance of event-driven and resource-aware design in sensor-network systems~\cite{tinyos}. These operating systems make embedded applications feasible, but they do not by themselves provide a cloud-side mechanism for expressing intent, binding identity, reconciling state, and verifying deployment outcomes.

WebAssembly adds a portable execution layer that can complement embedded operating systems. The original WebAssembly design emphasized safe, compact, and efficient portable code~\cite{haas_wasm,wasm_core}, while the WebAssembly System Interface (WASI) extended the ecosystem toward non-browser execution~\cite{wasi}. Runtime implementations such as WAMR, Wasmtime, WasmEdge, and wasm3 demonstrate that WASM can execute outside browsers and, in some cases, in resource-constrained settings~\cite{wamr,wasmtime,wasmedge,wasm3}. Container and cloud ecosystems have also begun to integrate WASM through projects such as runwasi and Spin~\cite{runwasi,spin}. These efforts establish runtime feasibility, but not operational orchestration: the unresolved question is how a cloud control plane can safely decide that a device should run a given module, transmit that intent through a trusted gateway, and receive evidence that execution and liveness match the requested state.

Secure IoT communication relies on standardized transport, object security, and compact serialization mechanisms. TLS 1.3 and DTLS 1.3 provide channel security~\cite{rfc8446,rfc9147}; CoAP and MQTT address lightweight application-layer communication~\cite{rfc7252,mqtt}; OSCORE and COSE provide object security and compact cryptographic containers~\cite{rfc8613,rfc9052}; and CBOR supports low-overhead binary serialization~\cite{rfc8949}. These mechanisms are complementary to the present work. A secure channel can protect messages in transit, but it does not determine which component owns reported state, when a device should be considered connected, or how a cloud resource should converge after a deployment acknowledgment.

\subsection{Reconciliation, Evaluation, and Positioning}
The operator pattern and Kubernetes controller model encode desired state, observe actual state, and repeatedly reconcile the two~\cite{kubernetes,kubernetes_operators}. In conventional clusters, this model works because nodes and workloads participate directly in the Kubernetes ecosystem. Edge systems complicate the model because connectivity can be intermittent, observations may be delayed, and devices may expose limited introspection. Prior platforms address parts of this challenge through edge agents, device twins, or local gateways~\cite{kubeedge,openyurt,akri,edgex,lwm2m}. For embedded targets, however, the gateway becomes more than a message forwarder: it is the semantic boundary between cloud resources and protocol-derived evidence.

Reproducible evaluation is equally important because physical hardware availability, network conditions, and manual setup can make embedded IoT experiments difficult to repeat. Large-scale experimental infrastructures such as SLICES highlight the role of programmable, shared scientific instruments for networking research~\cite{fdida2022slices}. Renode supports hardware-faithful emulation and multi-node embedded scenarios, while Cooja, ns-3, and OMNeT++ support repeatable simulation and network experimentation in wireless, sensor, and IoT research~\cite{renode,cooja,ns3,omnet}. The evaluation in this paper follows that direction by using an emulation-driven validation path and structured measurement records. This is appropriate because the main claim concerns end-to-end orchestration coherence rather than radio performance, energy consumption, or microarchitectural timing.

%Removed (ridondanza): Taken together, existing cloud-edge orchestration systems extend Kubernetes toward capable edge nodes; IoT middleware improves interoperability and telemetry; embedded operating systems and WASM runtimes make constrained execution feasible; and secure communication standards provide protected message exchange. The proposed framework connects these strands at a narrower integration point: a Kubernetes-native control plane expressing declarative deployment intent for an embedded-class endpoint that remains outside the cluster and reports execution evidence through a gateway.
Unlike middleware-centric approaches, the main abstraction is not a device model or telemetry bus, but the convergence between CRD intent, gateway evidence, and embedded WASM execution. The contribution is therefore best understood as a systems-integration study that combines explicit identity binding, gateway-mediated control, heartbeat-based liveness, CRD status convergence, and reproducible quantitative validation.

\begingroup
\color{red}
Table~\ref{tab:related-positioning} summarizes frameworks' characteristics. The comparison is qualitative because the surveyed systems target different layers of the edge continuum, but it clarifies the specific gap addressed here: preserving Kubernetes-native intent and reconciliation while keeping the embedded endpoint outside the cluster and still obtaining execution evidence.

\begin{table*}[pos=t]
\centering
\caption{Qualitative comparison relative to representative edge, IoT, and WASM orchestration approaches.}
\label{tab:related-positioning}
\footnotesize
\setlength{\tabcolsep}{3pt}
\begin{tabular}{p{0.17\linewidth} p{0.20\linewidth} p{0.20\linewidth} p{0.18\linewidth} p{0.17\linewidth}}
\toprule
\textbf{Approach} & \textbf{Primary abstraction} & \textbf{Assumed edge capability} & \textbf{Kubernetes integration} & \textbf{Gap for this work} \\
\midrule
KubeEdge/OpenYurt & Cloud-edge node and application management & Capable edge hosts or edge nodes with local agents & Extends Kubernetes control toward the edge & Too heavy for endpoints that cannot host cluster agents. \\
Akri/device-plugin style integration & Discovery and exposure of leaf devices & Kubernetes nodes near the device & Represents devices through cluster-attached resources & Does not make the constrained device itself a verifiable WASM execution target. \\
Krustlet/runwasi-style WASM execution & WASM workload as schedulable unit & Host able to participate in Kubernetes scheduling & Integrates WASM into Kubernetes workload execution & Targets worker-like execution environments rather than non-node embedded firmware. \\
\revC{Edge FaaS/deviceless serverless} & \revC{Function as deployment unit, with capability- and load-aware placement} & \revC{Linux single-board edge workers able to host a container-based function runtime} & \revC{Kubernetes- and Docker-based deployment of the function runtime} & \revC{Assumes an edge worker that can run a container runtime, not a non-node embedded endpoint.} \\
IoT middleware & Device model, telemetry, interoperability, and management & Gateway, server, or device-management infrastructure & May be bridged to cloud platforms but is not primarily a CRD reconciliation model & Provides broad IoT management, but not Kubernetes-native deployment intent to embedded WASM execution evidence. \\
This framework & CRD intent plus gateway-mediated evidence & Zephyr/WAMR endpoint with TLS/CBOR protocol support only & Uses CRDs, controllers, and status reconciliation without making the device a worker node & Focused on reproducible orchestration of constrained embedded WASM targets. \\
\bottomrule
\end{tabular}
\end{table*}
\endgroup
